% Vietnamese translation for OI Public Library
% Source-File: IOI中国国家候选队论文/2004/汪汀.ppt
\documentclass[11pt]{article}
\usepackage{oipl-vn}

\newcommand{\Oh}{\mathcal{O}}

\title{Các mở rộng của cây khung nhỏ nhất\\{\large Ghi chú theo slide}}
\author{Wang Ting}
\date{IOI 2004}

\begin{document}
\maketitle

\origtitle{Extensions of the minimum spanning tree problem}
\sourcepath{IOI 2004 / Wang Ting.ppt}

\section{Nền tảng}

Slide bắt đầu từ đồ thị \(G=(V,E,w)\) và hai thuật toán quen thuộc:
Prim với heap có thể đạt \(\Oh(|V|\log |V|+|E|)\), còn Kruskal đạt
\(\Oh(|E|\log |V|)\). Từ nền tảng này, bài trình bày xét các biến thể khi
cây khung phải thỏa thêm điều kiện.

\section{Cây khung có ràng buộc bậc}

Một biến thể nổi bật là ràng buộc bậc tại đỉnh \(v_0\): cần cây \(T\) sao cho
\(d_T(v_0)=k\). Các slide dùng phép trao đổi cạnh \(T+a-b\), trong đó thêm
cạnh ngoài cây \(a\) rồi bỏ cạnh \(b\) trên chu trình vừa tạo. Tập
\((+a,-b)\) biểu diễn các trao đổi hợp lệ.

Để chọn trao đổi tốt, ta gốc hóa cây tại \(v_0\) và duy trì
\(\mathit{Father}(v)\), \(\mathit{Best}(v)\): cạnh nặng nhất trên đường từ
\(v_0\) tới \(v\). Khi thêm cạnh nối hai thành phần, cạnh tốt nhất cần bỏ
có thể tìm nhanh trên đường trong cây.

\section{Cây khung nhỏ thứ hai}

Với mỗi cạnh ngoài cây \((v_1,v_2)\), thêm cạnh ấy vào MST tạo ra đúng một
chu trình. Nếu bỏ cạnh nặng nhất trên đường \(v_1\) đến \(v_2\), ta thu được
một cây khung khác tốt nhất cho lựa chọn đó. Tiền xử lý \(\mathit{Best}\) cho
các cặp đỉnh giúp xét mọi cạnh ngoài cây trong \(\Oh(|E|)\) sau bước chuẩn
bị.

\section{Thông điệp}

Các mở rộng MST thường vẫn xoay quanh tính chất cắt, chu trình và trao đổi
cạnh. Khi thêm ràng buộc, hãy tìm phép trao đổi nhỏ nhất làm thay đổi đúng
đại lượng cần kiểm soát, rồi dùng cây hiện tại để trả lời nhanh phần ``bỏ
cạnh nào''.

\end{document}
